% !TeX root = ../tesis.tex


\chapter{Microcito}

\label{section:apendix5}

En la Sección \ref{section:Res_Microcito} se discutieron los resultados obtenidos de $Q_{\text{ext}}$ y del patrón de esparcimiento para $\lambda=431$ nm, correspondiente a la banda de Soret de un microcito con CHCM de 28.7 g/dL. En esta sección se presentan los resultados para las longitudes de onda representativas de las bandas $\beta$ y $\alpha$. En la Fig.~\ref{fig:FF_541_micro_1} se muestra la sección transversal diferencial de esparcimiento $\dd{C_{\text{sca}}}/\dd{\Omega}$, así como su representación polar [Fig. \ref{subfig:FF_polar_541_micro}] para los cortes $\varphi=0^{\circ}$ y $\varphi=90^{\circ}$ para $\lambda=541$ nm. En la Fig. \ref{fig:FF_cartesian_541_health} se muestra la representación cartesiana de $\dd{C_{\text{sca}}}/\dd{\Omega}$ para los cortes $\varphi=0^{\circ}$ y $\varphi=90^{\circ}$ para $\lambda=541$ nm. Se observa que en todas las longitudes de onda el esparcimiento hacia adelante es el predominante. Además, al igual que con el esferocito, se encuentran diferencias entre los cortes $\varphi=0^{\circ}$ y $\varphi=90^{\circ}$ en la región de retroesparcimiento, donde la correlación aumenta a $r=0.99.$ 
	%
\begin{figure}[h!]
	\centering
	\sidesubfloat[First image]{\hspace{-0.5cm}{	\includegraphics[width=7.3cm]{4-Resultados/FFdistribution541_Micro.pdf}\label{subfig:FF_541_micro}}}\hspace{0.1cm}
	\sidesubfloat[Second image]{\hspace{-0.5cm}{
			\includegraphics[width=7.6cm]{4-Resultados/polarplot541FFComEsfMicro.pdf}\label{subfig:FF_polar_541_micro}}}\vspace{-0.1cm}
	\caption{\textbf{a)} Sección transversal diferencial de esparcimiento $\dd{C_{\text{sca}}}/\dd{\Omega}$ normalizada respecto a su valor máximo para una longitud de onda $\lambda=541$ nm en función del ángulo de esparcimiento $\theta$ y del ángulo $\varphi$ para un eritrocito sano. Se identifican dos planos $\varphi=0^{\circ}$ en azul claro y $\varphi=90^{\circ}$ en azul oscuro. \textbf{b)} Sección transversal diferencial de esparcimiento $\dd{C_{\text{sca}}}/\dd{\Omega}$ en decibeles en función de $\theta$ para dos cortes: $\varphi=0^{\circ}$ (puntos azules claro) y $\varphi=90^{\circ}$ (puntos azules oscuro). La línea continua que une los puntos es una guía visual.}
	\label{fig:FF_541_micro_1}
\end{figure}
%

	%
\begin{figure}[h!]
	\centering
			\includegraphics[width=14cm]{4-Resultados/linearplot541FFMicro.pdf}\vspace{-0.1cm}
	\caption{Sección transversal diferencial de esparcimiento $\dd{C_{\text{sca}}}/\dd{\Omega}$ en función de $\theta$ para dos cortes: $\varphi=0^{\circ}$ (puntos azules claro) y $\varphi=90^{\circ}$ (puntos azules oscuro). La línea continua que une los puntos es una guía visual.}
	\label{fig:FF_cartesian_541_micro}
\end{figure}
%
En la Fig.~\ref{subfig:FF_573_micro} se muestra la sección transversal diferencial de esparcimiento $\dd{C_{\text{sca}}}/\dd{\Omega}$, así como su representación polar [Fig. \ref{subfig:FF_polar_573_micro}] para los cortes $\varphi=0^{\circ}$ y $\varphi=90^{\circ}$ para $\lambda=573$~nm. En la Fig. \ref{fig:FF_cartesian_573_micro} se muestra la representación cartesiana de $\dd{C_{\text{sca}}}/\dd{\Omega}$ para los cortes $\varphi=0^{\circ}$ y $\varphi=90^{\circ}$ para $\lambda=573$ nm. Al igual que para las longitudes de onda de 431 y 541 nm, el esparcimiento se concentra predominantemente en la región frontal. En esta longitud de onda, disminuyen las diferencias entre los cortes $\varphi=0^{\circ}$ y $\varphi=90^{\circ}$ donde la correlación aumenta a $r=0.93$ en la región lateral y a $r=0.99$ en la región de retroesparcimiento. 
	%
\begin{figure}[h!]
	\centering
	\sidesubfloat[First image]{\hspace{-0.5cm}{	\includegraphics[width=7.3cm]{4-Resultados/FFdistribution573_Micro.pdf}\label{subfig:FF_573_micro}}}\hspace{0.1cm}
	\sidesubfloat[Second image]{\hspace{-0.5cm}{
			\includegraphics[width=7.6cm]{4-Resultados/polarplot573FFComEsfMicro.pdf}\label{subfig:FF_polar_573_micro}}}\vspace{-0.1cm}
	\
	\vspace{0.1cm}
	\caption{\textbf{a)} Sección transversal diferencial de esparcimiento $\dd{C_{\text{sca}}}/\dd{\Omega}$ normalizada respecto a su valor máximo para una longitud de onda $\lambda=573$ nm en función del ángulo de esparcimiento $\theta$ y del ángulo $\varphi$ para un eritrocito sano. Se identifican dos planos $\varphi=0^{\circ}$ en azul claro y $\varphi=90^{\circ}$ en azul oscuro. \textbf{b)} Sección transversal diferencial de esparcimiento $\dd{C_{\text{sca}}}/\dd{\Omega}$ en decibeles en función de $\theta$ para dos cortes: $\varphi=0^{\circ}$ (puntos azules claro) y $\varphi=90^{\circ}$ (puntos azules oscuro). La línea continua que une los puntos es una guía visual.}
	\label{fig:FF_573_micro_1}
\end{figure}
%

	%
\begin{figure}[h!]
	\centering
	
			\includegraphics[width=14cm]{4-Resultados/linearplot573FFMicro.pdf}\vspace{-0.1cm}
	\caption{Sección transversal diferencial de esparcimiento $\dd{C_{\text{sca}}}/\dd{\Omega}$ en función de $\theta$ para dos cortes: $\varphi=0^{\circ}$ (puntos azules claro) y $\varphi=90^{\circ}$ (puntos azules oscuro). La línea continua que une los puntos es una guía visual.}
	\label{fig:FF_cartesian_573_micro}
\end{figure}
%

Se observa que conforme aumenta la longitud de onda, la dirección de esparcimiento predomina más hacia adelante. Asimismo, se observa que la diferencia entre los cortes $\varphi=0^{\circ}$ y $\varphi=90^{\circ}$ dismimuye al incrementar la longitud de onda de la luz incidente.